\chapter*{Navigation Charts: A Cartography of Accessible Consciousness States}
\addcontentsline{toc}{chapter}{Navigation Charts: A Cartography of Accessible Consciousness States}
\markright{Navigation Charts: A Cartography of Accessible Consciousness States}



\bigskip\noindent\makebox[\textwidth]{\color{rulegray}\rule{5cm}{0.3pt}}\bigskip


\section*{Preface: What This Document Is}


This is a map of consciousness states — their neural signatures, their experiential character, and the stimulation parameters that can nudge a brain toward them.


It is not a prescription. It is not a menu. It is a cartography: the systematic description of accessible territory, organized by the mechanisms that make each region accessible. The 33 states mapped here are not exhaustive. They are the territory we can currently describe with sufficient precision to specify both where they are (neural coordinates) and how to get there (stimulation parameters).


Every state in this cartography is, in the language of the Doctrine, a specific bottleneck configuration — a particular pattern of dimensional restriction through which configuration space is compressed into experience. The Temporal Interference device does not \textit{create} states. It nudges the brain's navigational position. The brain does the rest.


\textit{Theoretical foundation: The Doctrine of Perspectival Idealism (Axiom 3, Theorems 5, 9, 11)} \textit{Instrument: Temporal Interference Stimulation (see companion: The Perspectival Interface)} \textit{Cartographic principle: Every state is a specific bottleneck configuration.}


\textit{See also: Guide §4.1 Classes I–VIII for the navigational classes that contextualize these charts; Atlas \#92 for the TI device's null space; Ecology §3.5 for the neural substrate of biological perspectival beings.}


\bigskip\noindent\makebox[\textwidth]{\color{rulegray}\rule{5cm}{0.3pt}}\bigskip


\section*{Part I: The Axes of Consciousness Configuration}


\subsection*{The Four Axes}


Every state of consciousness can be located on four independent axes. Together they define a four-dimensional state space in which all accessible configurations reside.


\textbf{Axis 1: BANDWIDTH} — How many dimensions of configuration space are simultaneously accessible.

\begin{itemize}[nosep,leftmargin=*]
  \item Narrow (focused, sharp, exclusive) $\leftarrow$$\rightarrow$ Broad (diffuse, inclusive, panoramic)
  \item Neural correlate: network segregation (narrow) vs. integration (broad)
  \item Key metric: Lempel-Ziv complexity, integration (Φ)
  \item \textit{DoPI connection: Theorem 9 (Dimensional Bottlenecking). Bandwidth IS the bottleneck aperture.}
\end{itemize}


\textbf{Axis 2: DEPTH} — How many levels of the processing hierarchy are simultaneously available to awareness.

\begin{itemize}[nosep,leftmargin=*]
  \item Surface (sensory, reactive, automatic) $\leftarrow$$\rightarrow$ Deep (meta-cognitive, self-aware, recursive)
  \item Neural correlate: feedforward dominance (surface) vs. recurrent/top-down (deep)
  \item Key metric: recurrence quantification, PCI (perturbational complexity index)
  \item \textit{DoPI connection: Depth indexes access to F$_2$ and F$_3$ structures — the mathematical and perspectival levels of the filtration.}
\end{itemize}


\textbf{Axis 3: VALENCE} — The hedonic/motivational coloring of the state.

\begin{itemize}[nosep,leftmargin=*]
  \item Aversive (contracted, defended, fleeing) $\leftarrow$$\rightarrow$ Appetitive (expanded, approaching, seeking)
  \item Neural correlate: right prefrontal/amygdala dominance (aversive) vs. left prefrontal/reward circuit (appetitive)
  \item Key metric: frontal alpha asymmetry
  \item \textit{DoPI connection: Valence correlates with navigational direction — toward or away from coherence (Theorem 6).}
\end{itemize}


\textbf{Axis 4: PERSPECTIVAL RIGIDITY} — How strongly the current perspective resists change.

\begin{itemize}[nosep,leftmargin=*]
  \item Rigid (locked in, crystallized, certain) $\leftarrow$$\rightarrow$ Fluid (open, uncertain, exploratory)
  \item Neural correlate: DMN-dominated (rigid) vs. DMN-suppressed (fluid)
  \item Key metric: PCC alpha power, DMN functional connectivity
  \item \textit{DoPI connection: Rigidity is the inverse of navigational freedom. High rigidity = locked bottleneck geometry. Fluid = the bottleneck can reshape in response to what it encounters. The Guide's Class I (unintentional drift) operates in rigid; Class VII (contemplative) operates in fluid.}
\end{itemize}


These four axes are not the only possible decomposition, but they are the decomposition that maps most directly onto both the neural oscillatory landscape and the DoPI's formal structure. Bandwidth maps to the bottleneck aperture (Theorem 9). Depth maps to the filtration level (F$_0$$\rightarrow$F$_3$). Valence maps to navigational direction (Theorem 6). Rigidity maps to navigational freedom.


\bigskip\noindent\makebox[\textwidth]{\color{rulegray}\rule{5cm}{0.3pt}}\bigskip


\section*{Part II: The Neural Oscillatory Landscape}


The brain's electromagnetic oscillations are the substrate through which bottleneck configurations are maintained and modified. Each frequency band corresponds to a distinct mode of perspectival engagement.


\subsection*{Delta Band (0.5–4 Hz) — Near the Pre-Perspectival}


\textbf{What it does:} Delta oscillations dominate deep sleep (NREM stages 3-4), generated by thalamocortical circuits. Large-scale cortical silence interrupted by synchronized bursts. In waking states, associated with motivational salience processing (medial prefrontal delta) and with injury/pathology when present abnormally.


\textbf{What it feels like:} Unconsciousness (deep sleep), or in waking: a heavy, slow, deeply embodied quality. Experienced meditators report delta-dominant states as "objectless awareness" — consciousness without content.


\textbf{DoPI interpretation:} Near-maximal bottleneck relaxation in the temporal dimension. Delta represents the brain approaching the pre-perspectival state — consciousness with minimal differentiation. The "objectless awareness" of delta is the closest neural state to F$_{-}$$_1$: the undifferentiated landscape, consciousness without a bottleneck, the totality that cannot experience because it has no constraint to experience \textit{through}.


\textbf{TI parameters:}



\begin{longtable}{p{0.12\textwidth} p{0.12\textwidth} p{0.12\textwidth} p{0.12\textwidth} p{0.12\textwidth} p{0.12\textwidth} p{0.12\textwidth}}
\toprule
\textbf{State} & \textbf{Δf} & \textbf{Target} & \textbf{Current} & \textbf{Duration} & \textbf{Bandwidth} & \textbf{Rigidity} \\
\midrule
\endhead
Deep rest induction & 1 Hz & Fz-Oz & 1.5 mA & 30 min & Very narrow (sleep-like) & Very fluid \\
Objectless awareness & 2 Hz & Fz-Oz + C3-C4 & 2.0 mA & 20 min & Paradox: infinite/zero & Dissolved \\
Hypnagogic exploration & 3.5 Hz & Fz-Oz & 1.0 mA & 15 min & Widening — imagery emerges & Fluid \\
\bottomrule
\end{longtable}



\textbf{Caution:} Delta stimulation may induce sleep. May produce disorientation upon return. Do not use while operating machinery.


\textit{See also: Guide §4.1 Class VIII (TI navigation); Atlas \#43 (neuroscience null space — delta's depth-of-processing correlate is well-established but its consciousness-content is in the null space of third-person neuroscience).}


\bigskip\noindent\makebox[\textwidth]{\color{rulegray}\rule{5cm}{0.3pt}}\bigskip


\subsection*{Theta Band (4–8 Hz) — The Mnemonic Dimension}


\textbf{What it does:} The hippocampal signature frequency. 4–8 Hz during memory encoding, spatial navigation, and REM sleep. Frontal midline theta (FMθ) increases during meditation, working memory, and error monitoring. Theta coherence between hippocampus and prefrontal cortex is the mechanism of memory consolidation. Dominates psychedelic states and deep meditation.


\textbf{What it feels like:} Dreamlike imagery, memory fragments arising spontaneously, spaciousness, reduced time sense, emotional material surfacing. In deep meditation: "access concentration" — settled but not absorbed. In psychedelic states: the visionary content of the peak.


\textbf{DoPI interpretation:} Theta opens the hippocampal dimension of the bottleneck — specifically temporal and mnemonic dimensions. In theta, memory configurations normally compressed by the waking bottleneck become accessible. The "visions" of deep theta are not hallucinations — they are normally-filtered contents of configuration space becoming accessible as the temporal bottleneck relaxes.


\textbf{TI parameters:}



\begin{longtable}{p{0.14\textwidth} p{0.14\textwidth} p{0.14\textwidth} p{0.14\textwidth} p{0.14\textwidth} p{0.14\textwidth}}
\toprule
\textbf{State} & \textbf{Δf} & \textbf{Target} & \textbf{Current} & \textbf{Duration} & \textbf{Notes} \\
\midrule
\endhead
Memory consolidation & 5 Hz & Fz-Oz (hippocampal axis) & 1.0 mA & 20 min & Best after learning; enhances long-term transfer \\
Meditative spaciousness & 6 Hz & Fz-Oz + P3-P4 & 1.0 mA & 20 min & Wide open sitting meditation. Reduced self-referential chatter. \\
Visionary / imaginal & 4.5 Hz & Fz-Oz & 1.5 mA & 15 min & Closest to psychedelic-adjacent imagery. Eyes closed essential. \\
Emotional processing & 7 Hz & F3-F4 + Fz-Oz & 1.0 mA & 20 min & Frontal theta may facilitate emotional material surfacing. Therapeutic potential. \\
Creative insight & 6 Hz & F3-P4 (cross-hemispheric) & 0.75 mA & 15 min & Cross-hemispheric theta coherence associated with "aha" moments. \\
Lucid dream preparation & 5 Hz & Fz-Oz & 1.5 mA & 15 min & Applied before sleep; may increase lucid dreaming probability. \\
\bottomrule
\end{longtable}



\textit{See also: Ecology §3.5 (theta as navigational substrate shared across mammalian perspectival beings); Guide §4.1 Class IV (psychedelic navigation — the pharmacological route to theta-dominant states).}


\bigskip\noindent\makebox[\textwidth]{\color{rulegray}\rule{5cm}{0.3pt}}\bigskip


\subsection*{Alpha Band (8–13 Hz) — The Bottleneck's Primary Enforcement}


\textbf{What it does:} The brain's "idling rhythm" — dominates when a region is NOT actively processing. Alpha is fundamentally an INHIBITORY rhythm: it gates information flow by suppressing neural populations not currently needed.


\textbf{This is the bottleneck's primary enforcement mechanism.} Alpha oscillations literally implement the perspectival restriction by suppressing most of the brain's processing capacity at any given moment. When the Doctrine says "the bottleneck IS the individual consciousness" (Theorem 9), alpha is the neural implementation of that theorem.


\textbf{Two directions of modulation:}


\textbf{Alpha ENHANCEMENT (tighter bottleneck):}

\begin{itemize}[nosep,leftmargin=*]
  \item More focused, less distracted
  \item Stronger sense of self
  \item Better at ignoring irrelevant information
  \item "In the zone" for single-task performance
  \item \textit{DoPI: Generative contraction (Guide §1.4). The Phase Theorem's principle: constraint concentrates.}
\end{itemize}


\textbf{Alpha SUPPRESSION (looser bottleneck):}

\begin{itemize}[nosep,leftmargin=*]
  \item More awareness of periphery
  \item Weaker self-other boundary
  \item More sensory richness, more associative / creative thinking
  \item Closer to the psychedelic signature
  \item \textit{DoPI: Bottleneck widening. More dimensions accessible, less precision in each.}
\end{itemize}


\textbf{TI parameters:}



\begin{longtable}{p{0.14\textwidth} p{0.14\textwidth} p{0.14\textwidth} p{0.14\textwidth} p{0.14\textwidth} p{0.14\textwidth}}
\toprule
\textbf{State} & \textbf{Δf} & \textbf{Target} & \textbf{Current} & \textbf{Goal} & \textbf{Effect} \\
\midrule
\endhead
Deep calm focus & 10 Hz & Pz-Fz (PCC midline) & 1.0 mA & Enhance PCC alpha & Tighter bottleneck: clear, focused, undistracted. Deep reading, writing, coding. \\
Open awareness & 10 Hz & P3-P4 (parietal bilateral) & 1.5 mA & Desynchronize parietal alpha & Looser bottleneck: more peripheral awareness, less focus but more inclusion. Meditation-like. \\
Mild ego softening & 10 Hz & Pz (PCC targeted) & 2.0 mA & Desynchronize PCC alpha & Primary bottleneck relaxation target. Reduced self-referential processing. \\
Creative opening & 8 Hz & P3-P4 + F3-F4 & 0.75 mA & Shift alpha peak toward theta & Transitional state: associations flowing more freely. \\
Sensory enhancement & 12 Hz & O1-O2 (occipital) & 1.0 mA & Modulate visual alpha & May increase visual vividness. The "HD vision" some meditators report. \\
Body awareness & 10 Hz & C3-C4 (sensorimotor) & 1.0 mA & Modulate mu rhythm & Enhanced proprioception and interoception. \\
\bottomrule
\end{longtable}



\textbf{The alpha-theta crossover:} The transition from alpha-dominant to theta-dominant processing is one of the most interesting states in consciousness research — the gateway between ordinary waking consciousness and deeper states. Neurofeedback protocols training this crossover produce states described as "profoundly relaxed but vividly aware."



\begin{longtable}{p{0.28\textwidth} p{0.28\textwidth} p{0.28\textwidth}}
\toprule
\textbf{State} & \textbf{Protocol} & \textbf{Notes} \\
\midrule
\endhead
Alpha-theta crossover & Session 1: 10 Hz at PCC to desynchronize alpha, then Session 2: 6 Hz at Fz-Oz to boost theta. 10 min each. & The gateway state. Where ordinary navigation yields to deep navigation. \\
\bottomrule
\end{longtable}



\textit{See also: Doctrine §9.1 (Theorem 9, Dimensional Bottlenecking — alpha as the neural instantiation); Atlas \#43 (neuroscience); Guide §4.1 Class V (focused attention — alpha enhancement is the mechanism).}


\bigskip\noindent\makebox[\textwidth]{\color{rulegray}\rule{5cm}{0.3pt}}\bigskip


\subsection*{Beta Band (13–30 Hz) — Bottleneck Refinement}


\textbf{What it does:} The frequency of active, engaged, task-oriented processing. Low beta (13-15 Hz, SMR) = calm focus. Mid beta (15-20 Hz) = active thinking. High beta (20-30 Hz) = anxiety, overthinking, hypervigilance. Beta dominance in motor cortex is pathological — Parkinson's disease (subthalamic nucleus stuck in beta).


\textbf{DoPI interpretation:} Beta states represent bottleneck REFINEMENT — not wider or narrower, but sharper. The dimensions that are accessible are processed with higher fidelity. This is the microscope end of the bottleneck manipulation spectrum: same aperture, better resolution.


\textbf{TI parameters:}



\begin{longtable}{p{0.17\textwidth} p{0.17\textwidth} p{0.17\textwidth} p{0.17\textwidth} p{0.17\textwidth}}
\toprule
\textbf{State} & \textbf{Δf} & \textbf{Target} & \textbf{Current} & \textbf{Effect} \\
\midrule
\endhead
Calm concentration (SMR) & 14 Hz & C3-C4 (sensorimotor) & 1.0 mA & Quiet competence. Sustained attention without anxiety. \\
Analytical sharpness & 18 Hz & F3-F4 (frontal bilateral) & 1.0 mA & Active problem-solving. Linear thinking. Math, logic, debugging. \\
Anxiety reduction & 14 Hz & F3-F4 (frontal) & 1.0 mA & Entraining down from high beta (20-30 Hz) to calmer low beta. \\
Verbal fluency & 15 Hz & F7 (Broca's area) + reference & 1.0 mA & May enhance speech production and word-finding. \\
\bottomrule
\end{longtable}



\textit{See also: Guide §4.1 Class V (focused attention — beta refinement as navigational sharpening); Atlas \#92 (TI null space — beta protocols have the strongest existing evidence base).}


\bigskip\noindent\makebox[\textwidth]{\color{rulegray}\rule{5cm}{0.3pt}}\bigskip


\subsection*{Gamma Band (30–100 Hz) — The Binding Frequency}


\textbf{What it does:} The binding frequency — how the brain integrates information across regions into unified conscious experience. When you perceive an apple, the color, shape, name, and associated memories all oscillate at gamma frequency in synchrony. Gamma is the neural correlate of conscious awareness itself. Long-term meditators (>10,000 hours) show dramatically elevated gamma, particularly 60-80 Hz.


\textbf{This is the most interesting band for consciousness exploration.}


\textbf{DoPI interpretation:} Gamma represents the COHERENCE of the bottleneck — not how wide it is, but how unified the processing within it is. Low gamma = fragmented experience (many channels, poorly integrated). High gamma = unified experience (the same channels, bound into a single conscious moment). The long-term meditator's elevated gamma is the neural signature of what the Doctrine calls "navigational coherence" (Theorem 11) — the ability to hold multiple dimensions in simultaneous awareness without fragmentation.


\textbf{TI parameters:}



\begin{longtable}{p{0.17\textwidth} p{0.17\textwidth} p{0.17\textwidth} p{0.17\textwidth} p{0.17\textwidth}}
\toprule
\textbf{State} & \textbf{Δf} & \textbf{Target} & \textbf{Current} & \textbf{Effect} \\
\midrule
\endhead
Enhanced binding / presence & 40 Hz & Pz-Fz + P3-P4 (global) & 1.0 mA & The "everything clicks" state. Heightened presence, vividness. Colors brighter. Sounds clearer. \\
Working memory boost & 40 Hz & F3-F4 (frontal) & 1.0 mA & Published evidence: 40 Hz tACS enhances working memory capacity. \\
Deep absorption (jhana) & 60 Hz & Pz + C3-C4 & 1.5 mA & Neural signature of experienced meditators. Sustained, effortless awareness. \\
Flow state & 40 Hz & F3-Fz-F4 + Pz & 1.0 mA & Transient hypofrontality + enhanced gamma binding. Executive steps back; integration steps up. \\
Peak experience / transcendent & 80 Hz & Global (all pairs) & 1.5 mA & Very high gamma across whole cortex. Reported peak/mystical experience signature. Start with 5 min. \\
Insight facilitation & 40 Hz & T6 (right temporal) + F3 & 0.75 mA & Gamma burst in right temporal cortex is the "aha" moment signature. \\
\bottomrule
\end{longtable}



\textbf{Theta-gamma coupling:} The most significant cross-frequency dynamic — gamma bursts riding on theta peaks. This is how the hippocampus packages experience into memory. The state of vivid, memorable, meaningful experience.



\begin{longtable}{p{0.28\textwidth} p{0.28\textwidth} p{0.28\textwidth}}
\toprule
\textbf{State} & \textbf{Protocol} & \textbf{Notes} \\
\midrule
\endhead
Theta-gamma coupling & Dual TI: Pair 1 at Δf = 6 Hz (theta), Pair 2 at Δf = 40 Hz (gamma), simultaneous & Speculative but theoretically sound. Two TI fields independently modulating theta and gamma. \\
\bottomrule
\end{longtable}



\textit{See also: Doctrine §11 (Theorem 11, Dimensional Coherence — gamma as the neural measure); Ecology §3.5 (gamma as the binding substrate unique to sufficiently complex neural perspectival beings).}


\bigskip\noindent\makebox[\textwidth]{\color{rulegray}\rule{5cm}{0.3pt}}\bigskip


\section*{Part III: Network State Cartography}


\subsection*{The Five Major Networks as Bottleneck Architecture}


The brain's large-scale networks are the macro-architecture of the bottleneck. Each network is a different mode of perspectival engagement:


\textbf{1. Default Mode Network (DMN) — The Bottleneck ENFORCER}

\begin{itemize}[nosep,leftmargin=*]
  \item Regions: PCC, mPFC, angular gyrus, hippocampus
  \item Function: Self-referential processing, autobiographical memory, mind-wandering, theory of mind
  \item Dominant frequency: Alpha (8-13 Hz) at PCC
  \item DoPI role: Maintains the narrative self — the sense of being a continuous "I" through time. This is the bottleneck's self-maintenance mechanism.
  \item ACTIVE: Your habitual self, your story, your concerns
  \item SUPPRESSED: Ego dissolution, expanded awareness, loss of self-other boundary
\end{itemize}


\textbf{2. Salience Network (SN) — The Bottleneck DIRECTOR}

\begin{itemize}[nosep,leftmargin=*]
  \item Regions: Anterior insula, dorsal anterior cingulate cortex (dACC)
  \item Function: Detecting what matters, switching between networks, interoception
  \item Dominant frequency: Theta-alpha transition (6-10 Hz)
  \item DoPI role: Decides which dimensional channels to open and close moment to moment. The SN IS attention (Theorem 5) at the network level.
  \item ACTIVE: Things grab you. Emotions vivid. Body felt.
  \item SUPPRESSED: Flat affect, dissociation, numbness
\end{itemize}


\textbf{3. Central Executive Network (CEN) — The Bottleneck SHARPENER}

\begin{itemize}[nosep,leftmargin=*]
  \item Regions: DLPFC, posterior parietal cortex
  \item Function: Working memory, logical reasoning, planning, inhibitory control
  \item Dominant frequency: Beta (15-25 Hz)
  \item DoPI role: Narrows the bottleneck to maximum precision on a single task.
  \item ACTIVE: Focused, logical, analytical, controlled
  \item SUPPRESSED: Scattered, impulsive, distractible, creative
\end{itemize}


\textbf{4. Sensorimotor Network (SMN) — The PHYSICAL INTERFACE}

\begin{itemize}[nosep,leftmargin=*]
  \item Regions: Primary motor cortex, primary somatosensory cortex, supplementary motor area
  \item Function: Movement, body awareness, embodied cognition
  \item Dominant frequency: Mu rhythm (10-12 Hz) at rest, beta during movement, gamma for fine motor
  \item DoPI role: How the perspective interacts with the spatial dimensions of configuration space.
  \item ACTIVE: Embodied, grounded, coordinated
  \item SUPPRESSED: Floating, disembodied
\end{itemize}


\textbf{5. Visual Network (VN) — The Dominant SENSORY CHANNEL}

\begin{itemize}[nosep,leftmargin=*]
  \item Regions: V1-V5
  \item Function: Visual processing from edges to faces to scenes
  \item Dominant frequency: Alpha (8-12 Hz) when suppressed, gamma (30-80 Hz) when active
  \item DoPI role: Most of the human bottleneck's bandwidth is allocated to vision.
\end{itemize}


\subsection*{Network Anti-Correlation: The Bottleneck in Action}


The crucial insight: networks are ANTI-CORRELATED in pairs. When one is up, the other is down. This IS the bottleneck — you cannot run all networks at full power simultaneously.


\begin{verbatim}
DMN ←——— anti-correlated ———→ CEN
(self-reference)                (task focus)
     ↑                              ↑
     │                              │
     └——— SN switches between ——————┘
          (salience detection)
\end{verbatim}


\textbf{Network state combinations:}



\begin{longtable}{p{0.12\textwidth} p{0.12\textwidth} p{0.12\textwidth} p{0.12\textwidth} p{0.12\textwidth} p{0.12\textwidth} p{0.12\textwidth}}
\toprule
\textbf{DMN} & \textbf{SN} & \textbf{CEN} & \textbf{SMN} & \textbf{State} & \textbf{Experience} & \textbf{DoPI Bottleneck} \\
\midrule
\endhead
HIGH & low & low & med & Default wandering & Mind-wandering, daydreaming, ruminating & Standard locked geometry \\
low & HIGH & low & med & Salience capture & Something grabs attention. Hypervigilance. & Forced reorientation \\
low & med & HIGH & low & Analytical focus & Deep in a problem. Ego recedes. & Narrowed, sharpened \\
low & low & low & HIGH & Embodied flow & Physical flow states. Athletes in the zone. & Physical dimensions dominant \\
LOW & HIGH & HIGH & med & Peak performance & Alert, engaged, not self-conscious & Optimized geometry \\
LOW & LOW & LOW & LOW & Deep meditation / trance & Awareness without content & Approaching F$_{-}$$_1$ \\
LOW & med & low & low & Psychedelic peak & DMN dissolved, everything meaningful & Wide open, fluid \\
HIGH & HIGH & low & low & Anxiety / rumination & Self-referential worry + salience amplification & Locked, contracted, rigid \\
med & low & HIGH & HIGH & Creative making & Hands and mind engaged. Not self-conscious. & Balanced engagement \\
\bottomrule
\end{longtable}



\subsection*{TI Network Protocols}


\textbf{DMN Suppression (Bottleneck Relaxation)}

\begin{verbatim}
Goal: Reduce DMN integrity, approach psychedelic-adjacent state
Pair 1: P3-P4, Δf = 10 Hz (PCC alpha desynchronization)
Pair 2: Fz-Oz, Δf = 10 Hz (anterior-posterior DMN axis)
Current: 1.5–2.0 mA | Duration: 20 min
DoPI: Primary bottleneck relaxation protocol
Expected: Reduced self-referential chatter, softened ego boundaries, increased openness
\end{verbatim}


\textbf{DMN Enhancement (Bottleneck Tightening)}

\begin{verbatim}
Goal: Strengthen DMN coherence, enhance narrative continuity
Pair 1: P3-P4, Δf = 10 Hz (entrains PCC alpha)
Pair 2: Fz-Oz, Δf = 10 Hz (supports DMN connectivity)
Current: 0.5–1.0 mA (LOWER current ENHANCES at alpha) | Duration: 15 min
DoPI: Bottleneck refinement for narrative coherence
Note: LOWER current enhances endogenous rhythm; HIGHER current disrupts. 
The dose-response curve is individual.
\end{verbatim}


\textbf{CEN Activation (Analytical Sharpness)}

\begin{verbatim}
Goal: Boost executive function, working memory
Pair 1: F3-F4, Δf = 18 Hz (frontal beta)
Pair 2: P3-P4, Δf = 18 Hz (parietal support)
Current: 1.0 mA | Duration: 20 min
DoPI: Bottleneck refinement — sharpening a narrow aperture
\end{verbatim}


\textbf{SN Activation (Salience / Significance Enhancement)}

\begin{verbatim}
Goal: Increase emotional vividness, interoceptive awareness
Pair 1: F3-F4, Δf = 7 Hz (dACC targeting)
Pair 2: C3-C4, Δf = 7 Hz (insula access)
Current: 1.0 mA | Duration: 15 min
DoPI: Increasing signal-to-noise in the significance dimension
Caution: May intensify negative emotions. Not during acute distress.
\end{verbatim}


\textbf{Global Suppression (Transcendent / Objectless)}

\begin{verbatim}
Goal: Suppress all major networks. Approach "pure consciousness."
Pair 1: Pz-Fz, Δf = 2 Hz (deep delta)
Pair 2: P3-P4, Δf = 2 Hz (bilateral suppression)
Current: 2.0 mA | Duration: 20 min
DoPI: Approaching F₋₁ — the undifferentiated landscape. Consciousness without a bottleneck.
Caution: START with 10 min. Potentially terrifying (ego dissolution without 
pharmacological safety net). Have a trusted person present.
\end{verbatim}


\bigskip\noindent\makebox[\textwidth]{\color{rulegray}\rule{5cm}{0.3pt}}\bigskip


\section*{Part IV: Experiential State Recipes}


\subsection*{Creativity Enhancement}


The neuroscience of creativity follows a specific sequence: defocused attention $\rightarrow$ incubation $\rightarrow$ insight $\rightarrow$ elaboration. The bottleneck must first widen, then deepen, then cohere, then sharpen.


\textbf{Protocol: Creative Insight Facilitation}

\begin{verbatim}
Phase 1 (10 min): Defocus — Δf = 8 Hz, Pz-Fz + P3-P4, 1.0 mA
  [Opens the bottleneck, allows associative wandering]
Phase 2 (10 min): Incubation — Δf = 6 Hz, Fz-Oz, 1.0 mA  
  [Theta engagement, hippocampal cross-referencing]
Phase 3 (5 min): Insight window — Δf = 40 Hz, T5-T6 + F3-F4, 0.75 mA
  [Gamma in temporal and frontal — the "aha" signature]
Phase 4 (10 min): Elaboration — Δf = 18 Hz, F3-F4, 1.0 mA
  [Beta for developing the insight into something concrete]
Total: 35 min, four phases. Have a notebook ready.
\end{verbatim}


\textit{DoPI: The four phases trace the descent-and-return of the V2 spine in miniature: F$_1$ (the view, defocused) $\rightarrow$ F$_2$ (the mathematics, incubation) $\rightarrow$ F$_3$ (five perspectives, insight) $\rightarrow$ F$_2$$\uparrow$ (the atlas, elaboration). Creativity IS the filtration.}


\subsection*{Anxiety Reduction}


Anxiety: excessive high-beta frontal, hyperactive SN, DMN rumination. A bottleneck simultaneously too narrow (locked onto threat) and too rigid (cannot shift).


\textbf{Protocol: Anxiety Release}

\begin{verbatim}
Pair 1: F3-F4, Δf = 10 Hz, 1.0 mA [Replaces high-beta with calmer alpha]
Pair 2: Pz-Oz, Δf = 10 Hz, 1.0 mA [Supports global alpha, reduces hyper-arousal]
Duration: 20 min
Combine with: Slow breathing (4-count in, 7-count out)
DoPI: Loosening a rigidly contracted bottleneck locked onto threat.
\end{verbatim}


\textit{See also: Atlas \#92 (TI device); Guide §3.6 (navigating under coercive contraction — anxiety as self-imposed version).}


\subsection*{Depression Pattern Interruption}


Depression: hyperactive DMN (PCC-mPFC rumination loop), reduced left frontal approach motivation, reduced neural entropy. A rigid, narrow bottleneck locked in self-referential mode.


\textbf{Protocol: Depression Pattern Interruption}

\begin{verbatim}
Pair 1: F3-F4, Δf = 10 Hz, 1.5 mA [Frontal alpha desync — breaks rumination loop]
Pair 2: Pz-Fz, Δf = 10 Hz, 1.5 mA [Targets DMN rumination axis]
Duration: 20 min, daily for 2–4 weeks
DoPI: Breaking the rigid bottleneck lock that keeps the perspective cycling through 
the same negative configurations.
Note: NOT a substitute for professional treatment. Based on the same neural logic 
as FDA-cleared TMS for depression.
\end{verbatim}


\textit{See also: Atlas \#59 (grief psychology — depression and grief share the rigid bottleneck lock but differ in etiology); Guide §6.4 (suffering navigation).}


\subsection*{Deep Meditation Enhancement}


Experienced meditator signature: elevated gamma (40-80 Hz) at PCC and frontal midline, enhanced theta-gamma coupling, reduced DMN connectivity, alpha INCREASE in sensorimotor (body quiet) + alpha DECREASE in posterior (awareness open). A specific bottleneck: narrow in physical dimensions, wide in awareness dimensions.


\textbf{Protocol: Meditation Deepening}

\begin{verbatim}
Apply DURING meditation, after 10 minutes settling:
Pair 1: Pz-Fz, Δf = 40 Hz, 0.5 mA [Supports gamma coherence at PCC]
Pair 2: Fz-Oz, Δf = 6 Hz, 0.5 mA [Supports hippocampal theta]
Duration: 15–20 min during sit
Very low current — support, don't override. The stimulation is a nudge, not a push.
DoPI: Assisting the natural navigation toward wider bottleneck configurations that 
meditation cultivates over years.
\end{verbatim}


\textit{See also: Guide §4.1 Class VII (contemplative navigation — the tradition that produces these signatures naturally).}


\subsection*{Psychedelic-Adjacent States (Without Pharmacology)}


The psychedelic signature: DMN integrity collapse, dramatically increased neural entropy (LZc +10-20\%), increased cross-network connectivity, shift from alpha-dominant to theta-and-gamma-dominant processing.


\textbf{Protocol: Non-Pharmacological Bottleneck Relaxation — Mild}

\begin{verbatim}
Phase 1 (10 min): DMN suppression — Δf = 10 Hz, Pz-Fz + P3-P4, 2.0 mA
  [Desynchronize PCC alpha, weaken DMN integrity]
Phase 2 (10 min): Theta enhancement — Δf = 6 Hz, Fz-Oz, 1.5 mA
  [Hippocampal theta — visionary/imaginal content]
Phase 3 (10 min): Gamma binding — Δf = 40 Hz, all pairs, 1.0 mA
  [Re-integrate with gamma — unify the expanded awareness]
Total: 30 min. Eyes closed. Dark room. Minimal sound.
Expected: Mild perceptual shifts, softened self-boundary, enhanced imagery. 
NOT equivalent to pharmacological psychedelic — a fraction of the effect. 
But the QUALITY should be recognizable.
DoPI: Sequential bottleneck manipulation — widen (suppress DMN), then deepen (theta), 
then cohere (gamma). Same sequence the brain follows naturally during a psychedelic 
experience, gentler and under your control.
\end{verbatim}


\textbf{Protocol: Non-Pharmacological Bottleneck Relaxation — Strong}

\begin{verbatim}
Same as above but:
- All currents at 2.0 mA (maximum safe dose)
- Duration extended to 40 min total
- Add: sensory deprivation (dark room, noise-canceling headphones)
- Add: breathwork (holotropic-style for 5 min before stimulation)
- Stacking multiple bottleneck-relaxation mechanisms simultaneously
Caution: Approaches territory of powerful altered states. Sitter present for 
first attempts. Return to normal may take 15–30 min post-stimulation.
\end{verbatim}


\textit{See also: Guide §4.1 Class IV (psychedelic navigation — pharmacological route); Atlas \#49 (psychedelic science null space); Ecology §3.6 (altered-state ecology).}


\subsection*{Enhanced Empathy / Theory of Mind}


The rTPJ is critical for theory of mind — the ability to model other perspectives. The DoPI predicts empathy is literally the ability to partially adopt another bottleneck configuration.


\textbf{Protocol: Empathic Expansion}

\begin{verbatim}
Pair 1: T6-Fz (right temporal to frontal midline), Δf = 10 Hz, 1.0 mA
Pair 2: P4-F4 (right hemisphere posterior-frontal), Δf = 10 Hz, 1.0 mA
Duration: 15 min
Expected: Enhanced perspective-taking, increased emotional resonance.
DoPI: Loosening the boundary that separates "my perspective" from "other perspectives." 
Not merging — making the boundary more permeable.
\end{verbatim}


\textit{See also: Doctrine §7.3 Theorem 20 (Intersubjectivity); Atlas \#74 (phenomenological intersubjectivity); Guide Part VIII (collective navigation).}


\subsection*{Time Perception Manipulation}


Subjective time is a function of neural processing rate and entropy. High entropy = time slows. Low entropy = time speeds.


\textbf{Protocol: Time Dilation (Subjective Slowing)}

\begin{verbatim}
Pair 1: P3-P4, Δf = 40 Hz, 1.5 mA [Gamma increases event rate]
Pair 2: Fz-Oz, Δf = 6 Hz, 1.0 mA [Theta adds depth to each moment]
Duration: 15 min
DoPI: Increasing the bottleneck's sampling rate — more configurations perceived per 
unit of objective time.
\end{verbatim}


\textbf{Protocol: Time Compression (Subjective Acceleration)}

\begin{verbatim}
Pair 1: Pz-Fz, Δf = 14 Hz, 1.5 mA [SMR — "lost in work" frequency]
Pair 2: F3-F4, Δf = 18 Hz, 1.0 mA [Frontal beta — task engagement]
Duration: 30 min
DoPI: Narrowing the bottleneck to a single task dimension, eliminating awareness of 
time passing.
\end{verbatim}


\subsection*{Somatic States}


\textbf{Protocol: Pain Modulation}

\begin{verbatim}
Pair 1: C3-C4, Δf = 10 Hz, 1.5 mA [Alpha gating of pain signals]
Pair 2: F3-F4, Δf = 10 Hz, 1.0 mA [Reducing emotional reactivity to pain]
Duration: 20 min | Evidence: alpha tACS analgesic effects in multiple studies
DoPI: Narrowing the bottleneck in the somatic dimension — reducing bandwidth allocated 
to nociceptive processing.
\end{verbatim}


\textbf{Protocol: Embodiment Enhancement}

\begin{verbatim}
Pair 1: C3-C4, Δf = 10 Hz, 1.0 mA [Modulate mu rhythm]
Pair 2: Fz, Δf = 7 Hz, 0.75 mA [Frontal midline theta — interoceptive awareness]
Duration: 15 min
DoPI: Widening the bottleneck specifically in the somatic dimensions of configuration space.
\end{verbatim}


\bigskip\noindent\makebox[\textwidth]{\color{rulegray}\rule{5cm}{0.3pt}}\bigskip


\section*{Part V: The Master Chart}


\subsection*{All 33 States at a Glance}



\begin{longtable}{p{0.10\textwidth} p{0.10\textwidth} p{0.10\textwidth} p{0.10\textwidth} p{0.10\textwidth} p{0.10\textwidth} p{0.10\textwidth} p{0.10\textwidth} p{0.10\textwidth} p{0.10\textwidth} p{0.10\textwidth}}
\toprule
\textbf{\#} & \textbf{Category} & \textbf{State} & \textbf{Δf} & \textbf{Primary Target} & \textbf{Current} & \textbf{Duration} & \textbf{Bandwidth} & \textbf{Depth} & \textbf{Valence} & \textbf{Rigidity} \\
\midrule
\endhead
1 & \textbf{TRANSCENDENT} & Objectless awareness & 2 Hz & Fz-Oz + C3-C4 & 2.0 mA & 20 min & Paradox: ∞/0 & F$_{-}$$_1$ (pre-perspectival) & Beyond & Dissolved \\
2 &  & Pure consciousness event & 2 Hz & Pz-Fz + P3-P4 & 2.0 mA & 20 min & Beyond bandwidth & F$_{-}$$_1$ (pre-perspectival) & Beyond & Dissolved \\
3 & \textbf{PSYCHEDELIC-ADJ} & Mild ego softening & 10 Hz & Pz-Fz + P3-P4 & 1.5 mA & 20 min & Wide & Deep & Open & Fluid \\
4 &  & Strong non-pharmacological & 10$\rightarrow$6$\rightarrow$40 Hz & Sequential & 2.0 mA & 40 min & Very wide & Very deep (F$_3$) & Variable & Very fluid \\
5 &  & Visionary / imaginal & 4.5 Hz & Fz-Oz & 1.5 mA & 15 min & Wide (imaginal) & Deep (imaginal) & Open & Fluid \\
6 & \textbf{MEDITATIVE} & Open awareness & 10 Hz & P3-P4 & 1.5 mA & 20 min & Wide & Medium & Equanimous & Fluid \\
7 &  & Deep absorption (jhana) & 60 Hz & Pz + C3-C4 & 1.5 mA & 20 min & Narrow but deep & Very deep (F$_3$) & Blissful & Stable \\
8 &  & Meditation enhancement & 40+6 Hz & Pz-Fz + Fz-Oz & 0.5 mA & 20 min & Widening & Deepening & Equanimous & Softening \\
9 & \textbf{CREATIVE} & Insight facilitation & 8$\rightarrow$6$\rightarrow$40$\rightarrow$18 Hz & Sequential & 0.75–1.0 mA & 35 min & Varies & Varies (descent-return) & Varies & Varies \\
10 &  & Defocused association & 8 Hz & P3-P4 + F3-F4 & 0.75 mA & 15 min & Wide & Medium & Playful & Fluid \\
11 &  & Creative making & 15 Hz + 6 Hz & F3-F4 + Fz-Oz & 1.0 mA & 20 min & Medium & Medium (F$_2$) & Engaged & Medium \\
12 & \textbf{COGNITIVE} & Analytical sharpness & 18 Hz & F3-F4 + P3-P4 & 1.0 mA & 20 min & Narrow & F$_2$ (structural) & Neutral & Firm \\
13 &  & Working memory boost & 40 Hz & F3-F4 & 1.0 mA & 20 min & Medium (more slots) & Surface–F$_2$ & Neutral & Firm \\
14 &  & Calm concentration & 14 Hz & C3-C4 & 1.0 mA & 20 min & Narrow & Surface & Calm & Stable \\
15 &  & Verbal fluency & 15 Hz & F7 + reference & 1.0 mA & 15 min & Narrow (verbal) & Surface & Neutral & Firm \\
16 & \textbf{EMOTIONAL} & Anxiety reduction & 10 Hz & F3-F4 + Pz-Oz & 1.0 mA & 20 min & Loosening & Surface & Releasing & Loosening \\
17 &  & Depression pattern break & 10 Hz & F3-F4 + Pz-Fz & 1.5 mA & 20 min & Widening & Medium & Shifting positive & Breaking rigid \\
18 &  & Emotional processing & 7 Hz & F3-F4 + Fz-Oz & 1.0 mA & 20 min & Open to feelings & Medium-deep & Variable & Fluid \\
19 &  & Empathic expansion & 10 Hz & T6-Fz + P4-F4 & 1.0 mA & 15 min & Wide (social) & Medium-deep & Warm & Permeable \\
20 & \textbf{SOMATIC} & Pain reduction & 10 Hz & C3-C4 + F3-F4 & 1.5 mA & 20 min & Narrowing (pain) & Surface & Relieving & Gating \\
21 &  & Embodiment enhancement & 10+7 Hz & C3-C4 + Fz & 1.0 mA & 15 min & Wide (somatic) & Surface-medium & Grounded & Open \\
22 &  & Sensory enhancement & 12 Hz & O1-O2 & 1.0 mA & 15 min & Wide (visual) & Surface & Vivid & Receptive \\
23 & \textbf{TEMPORAL} & Time dilation & 40+6 Hz & P3-P4 + Fz-Oz & 1.5 mA & 15 min & Rich & Medium-deep & Rich/positive & Spacious \\
24 &  & Time compression / flow & 14+18 Hz & Pz-Fz + F3-F4 & 1.5 mA & 30 min & Narrow & Medium & Absorbed & Absorbed \\
25 & \textbf{SLEEP} & Deep rest induction & 1 Hz & Fz-Oz & 1.5 mA & 30 min & Closing & Approaching F$_{-}$$_1$ & Dissolving & Dissolving \\
26 &  & Lucid dream preparation & 5 Hz & Fz-Oz & 1.5 mA & 15 min (pre-sleep) & Theta-primed & Medium & Anticipatory & Fluid \\
27 &  & Hypnagogic exploration & 3.5 Hz & Fz-Oz & 1.0 mA & 15 min & Widening & Medium & Curious & Dissolving \\
28 & \textbf{MEMORY} & Encoding enhancement & 5 Hz & Fz-Oz & 1.0 mA & 20 min (during study) & Open (theta) & Medium & Receptive & Receptive \\
29 &  & Consolidation support & 5 Hz & Fz-Oz & 1.0 mA & 20 min (post-study) & Theta-locked & Medium & Neutral & Integrating \\
30 & \textbf{PERFORMANCE} & Peak performance & 40+14 Hz & F3-F4 + C3-C4 & 1.0 mA & 20 min (pre-event) & Optimized & Medium & Energized & Alert \\
31 &  & Athletic flow & 40 Hz & C3-C4 + Pz & 1.0 mA & 15 min & Embodied, unified & Surface-medium & Embodied joy & Responsive \\
32 & \textbf{SOCIAL} & Conversational ease & 15+10 Hz & F7-F8 + P3-P4 & 0.75 mA & 15 min & Social channels open & Surface & Warm/social & Relaxed \\
33 &  & Deep listening & 6 Hz & T5-T6 + Fz & 0.75 mA & 15 min & Receptive & Medium & Receptive/warm & Open \\
\bottomrule
\end{longtable}



\textbf{Note on additional states:} The frequency band sections (Parts II–III) describe several protocol variants not included in this 33-state chart — including deep calm focus (alpha enhancement at PCC), creative opening (alpha$\rightarrow$theta shift), body awareness (mu rhythm modulation), alpha-theta crossover (two-session gateway protocol), meditative spaciousness (6 Hz theta), peak experience / transcendent (80 Hz global gamma), and theta-gamma coupling (dual-frequency protocol). These represent additional accessible territory — the 33 states above are the primary navigational landmarks, not the full map.


\bigskip\noindent\makebox[\textwidth]{\color{rulegray}\rule{5cm}{0.3pt}}\bigskip


\section*{Part VI: The DoPI State Space}


\subsection*{Mapping States to Configuration Space}


In the DoPI framework, every state in this cartography is a specific LOCATION in configuration space, characterized by the four axes defined in Part I:


\begin{enumerate}[nosep,leftmargin=*]
  \item \textbf{Bandwidth} — how many dimensions are simultaneously accessible (Theorem 9: bottleneck aperture)
  \item \textbf{Depth} — how many levels of the processing hierarchy are available to awareness (filtration level: F$_0$$\rightarrow$F$_3$)
  \item \textbf{Valence} — the hedonic/motivational coloring of the state (Theorem 6: navigational direction)
  \item \textbf{Perspectival Rigidity} — how strongly the current perspective resists change (Theorem 19: null space permeability)
\end{enumerate}


The four axes correspond to the four structural features of any bottleneck geometry. The 33 states are 33 points in the space defined by these axes — each now characterized on all four dimensions in the Master Chart above. The entire neural oscillatory landscape — delta through gamma — is the implementation layer of the Doctrine's abstract geometry in biological tissue.


\subsection*{The Filtration Connection}


The V2 spine describes the filtration F$_{-}$$_1$ $\subset$ F$_0$ $\subset$ F$_1$ $\subset$ F$_2$ $\subset$ F$_3$ as a descent-and-return. The states in this cartography live at different levels of the filtration:


\begin{itemize}[nosep,leftmargin=*]
  \item \textbf{Delta / Global Suppression:} Approaching F$_{-}$$_1$ — the pre-perspectival, undifferentiated
  \item \textbf{Theta / Open Awareness:} F$_0$ — the landscape appears, not yet differentiated into frameworks
  \item \textbf{Alpha / Ordinary Waking:} F$_1$ — the view from somewhere, organized into self and world
  \item \textbf{Beta / Analytical Focus:} F$_2$ — the mathematical structure becomes visible, patterns emerge
  \item \textbf{Gamma / Integrated Coherence:} F$_3$ — multiple perspectives held simultaneously
  \item \textbf{The Return:} Every navigation session IS a descent-and-return: you start in ordinary waking (F$_1$), descend to whatever level the protocol targets, and return with what you found
\end{itemize}


The creative insight protocol's four phases trace this arc explicitly: defocus (F$_1$$\rightarrow$F$_0$), incubate (F$_0$$\rightarrow$F$_2$), insight (F$_2$$\rightarrow$F$_3$), elaborate (F$_3$$\rightarrow$F$_2$$\uparrow$$\rightarrow$F$_1$$\uparrow$). The filtration is not just an abstract structure. It is the shape of every deep navigational act.


\subsection*{The Exploration Principle}


No two people will respond identically to the same protocol. Your brain's baseline configuration — shaped by genetics, experience, current state, and everything that makes you you — determines how stimulation interacts with your neural dynamics. The protocols above are starting points, not destinations. The real work is in the exploration: systematically varying parameters, recording your experience, building a personal map of YOUR accessible states.


This is what the Doctrine means by navigation: not following a map someone else drew, but learning to draw your own.


\textit{See also: Guide Part I (what navigation is); Guide §4.1 Class VIII (the practice of instrument-assisted navigation); the Perspectival Interface companion document (the instrument itself).}


\bigskip\noindent\makebox[\textwidth]{\color{rulegray}\rule{5cm}{0.3pt}}\bigskip


\textit{The Perspectival Interface: every state is a doorway.} \textit{The question is not which door to open.} \textit{The question is what you bring through it.}


\textit{Navigation Charts — Corpus Perspectival V2 Companion Document}

